\section*{1. 范围}\addcontentsline{toc}{section}{1. 范围}

{\normalsize
\subsection*{1.1 标识}\addcontentsline{toc}{subsection}{1.1 标识}

{\normalsize
（1）文档标识号：xxxxx\_STR\_V0.0.0.1；

（2）系统标识xxxxx，在测试文档中简称为xxxxx；

（3）项目名称：xxxxx系统以下简称xxxxx；

（4）文档名称：xxxxx节点软件测试报告；

（5）软件版本号：releasexxxxx1版；

（6）本文档适用于xxxxx系统2025年9月节点软件测试过程。
}

\subsection*{1.2 系统概述}\addcontentsline{toc}{subsection}{1.2 系统概述}

{\normalsize
系统用途部分内容敏感，不在此处列出。详见xxxxx系统合同（xxxxx）。

xxxxx规定的本阶段主要要求和技术指标为：

xxxxx。

项目节点测试的测试项与xxxxx的覆盖性对应关系如表\ref{tbl:coverage-map}所示。

{\settablespacing
\begin{longtblr}[theme=gjb,caption={xxxxx与测试项覆盖性对照表},label={tbl:coverage-map}]{
  colspec={|p{0.8cm}|p{6.8cm}|p{6.9cm}|},
  rowhead=1,
  hlines={wd=\GjbTableRuleWd,fg=\GjbTableRuleColor},
  column{1}={halign=c},
}
序号 & xxxxx & 测试项 \\
\Seq & xxxxx &  \\
\Seq & xxxxx &  \\
\Seq & xxxxx &  \\
\Seq & xxxxx &  \\
\Seq & xxxxx &  \\
\end{longtblr}
}
\vspace{-6pt}

xxxxx系统软件的需方是xxxxx管理办公室，开发方是xxxxx。

标识当前和计划运行的现场，测试地点为“xxxxx”，测试环境包括1台国产ARM架构服务器，其搭载国产处理器（飞腾）和国产操作系统（银河麒麟服务器版）以及2台国产X86台式机，其中服务器部署xxxxxreleasexxxxx1版，两台台式机为。
}

\subsection*{1.3 文档概述}\addcontentsline{toc}{subsection}{1.3 文档概述}

{\normalsize
本文档用于xxxxx节点测试报告，规定了本次测试的环境以及具体需要执行的测试用例，指导测试人员进行具体的测试活动。

本文档模板涵盖了GJB438C-2021《军用软件开发文档通用要求》对软件测试报告文档的要素和内容的要求。

描述本文档密级为xxxxx，必须按照xxxxx文件的要求使用和分发。即使在保密性得到保证的情况下，对本文档的使用也必须经研制方主管人员的认可。
}

}

% =================== 第2章：引用文档 ===================
\noindent
